Variáveis $(n=4)$: $Q, d, \mu, p/l$

Dimensões $(m=3)$: $M, L, T$

Número de termos de $\pi$: $r = 4 - 3 = 1$

\[ \pi_1 = Q d^\alpha \mu^\beta (p/l)^\gamma \]

Substituindo pelas respectivas dimensões:
\[ M^{0}L^{0}T^{0} = (L^{3}T^{-1}) \cdot (L)^{\alpha} \cdot (M L^{-1}T^{-1})^{\beta} \cdot (M L^{-2}T^{-2})^\gamma \]

Analisando cada dimensão:
\begin{align*}
\text{M:} \quad & 0 = \beta + \gamma \\
\text{T:} \quad & 0 = -1 - \beta - 2\gamma \\
\text{L:} \quad & 0 = 3 + \alpha - \beta - 2\gamma
\end{align*}

Resolvendo o sistema: $\gamma = -1$, $\beta = 1$ e $\alpha = -4$.

Substituindo os valores encontrados:
\[ \pi_1 = \frac{Q\mu}{d^4 (p/l)} \]

Rearranjando para encontrar a vazão $Q$:
\[ Q = (\text{constante}) \frac{d^{4}}{\mu(p/l)} \]